\begin{quotation} [1]
"The figures are helpful for showing how the existing tools work, but for the
proposal make sure you show how you want the visualized code to work.  This
could be through making a mockup or setting up input for Manim to generate
something that looks like what you want.  You should also think about the
pipeline for what data you are planning on getting from LLVM, how you are going
to get it, and how you will feed it into Manim."
\end{quotation}

\noindent This feedback remains unaddressed at present, primarily due to time constraints during the proposal development phase. We intend to create a mockup of the intended output in the near future to better illustrate our vision for the final product.

\begin{quotation} [2]
  "You have figures for some existing tools.  Make sure you talk about them in the proposal and what their limitations are."
\end{quotation}

\noindent\hrulefill

\begin{quotation} [3]
  "This [PythonTutor] figure is unreadable."
\end{quotation}

\begin{quotation} [4]
  "[PythonTutor] Spelling"
\end{quotation}

\begin{quotation} [5]
  "Which languages [are supported by PythonTutor]?"
\end{quotation}

\begin{quotation} [6]
  "You should have a citation for where to access the [PythonTutor] tool itself and not just the paper describing its efficacy."
\end{quotation}

\begin{quotation} [7]
  "The text in the [control-flow graph] figure is too small."
\end{quotation}

\begin{quotation} [8]
  "[Graphviz] Linky please"
\end{quotation}

\begin{quotation} [9]
  "Linky to Manim."
\end{quotation}

\begin{quotation} [10]
  "citation needed [for 3Blue1Brown's winning a Streamy award]"
\end{quotation}

\begin{quotation} [11]
  "The order of the [model] diagram is in the opposite order that they are presented in [implementation plan] text."
\end{quotation}

\begin{quotation} [12]
  "The items in the timeline are incredibly brief and need more detail or support."
\end{quotation}

\noindent\hrulefill

\begin{quotation} [13]
  "The figures are quite shrunk down especially figure 1 and figure 2 making it extremely difficult to read the content within it, if those were made larger it would greatly improve the readability."
\end{quotation}

\begin{quotation} [14]
  "The abstract clearly explains the purpose and motivation of LLVManim, especially its goal of improving visualization aesthetics for presentations and education. However, it could briefly mention the intended users explicitly, such as educators, students, or software engineers. Additionally, the abstract could briefly summarize how LLVManim achieves this goal rather than only stating the objective."
\end{quotation}

\begin{quotation} [15]
  "The introduction provides strong motivation and clearly explains the gap between current visualization tools and presentation-quality output. The problem statement is well articulated, especially the distinction between “function over form” tools and presentation needs. However, the introduction could more explicitly define the scope of the LLVManim system itself earlier in the section. Currently, LLVManim is only briefly mentioned at the end. Introducing LLVManim sooner would improve clarity and help frame the remainder of the document."
\end{quotation}

\begin{quotation} [16]
  "[The PythonTutor] section clearly explains PythonTutor’s strengths and limitations. The educational benefits and usability are well supported by references, which strengthens the argument. However, this section could be slightly condensed. Some explanations regarding educational benefits are repeated, and the limitations relevant specifically to LLVManim’s goals like their lack of customization, presentation-quality output, etc., should be emphasized more clearly. Additionally, explicitly comparing PythonTutor directly to LLVManim in one or two sentences would make the relevance stronger."
\end{quotation}

\begin{quotation} [17]
  "[The LLVM Analysis and Transform Passes] section provides an excellent explanation of LLVM’s capabilities and limitations. The explanation of LLVM IR, analysis passes, and DOT visualizations is clear and technically accurate. However, this section could be improved by separating the strengths and weaknesses into clearly defined subsections or paragraphs. Currently, strengths and limitations are interwoven, which makes it slightly harder to skim. Additionally, some technical explanations like the IR syntax examples could be shortened slightly, as the focus of the proposal is visualization rather than compiler internals."
\end{quotation}

\begin{quotation} [18]
  "[The Manim] section provides an excellent explanation of Manim’s capabilities and its usefulness for creating high-quality educational animations. The explanation of its role in improving visualization aesthetics is clear and well justified. However, this section could more explicitly explain how Manim integrates into the LLVManim system pipeline. Additionally, some historical background about 3Blue1Brown could be slightly condensed to maintain focus on the project itself."
\end{quotation}

\begin{quotation} [19]
  "[The Goals and Features] section clearly defines the goals, proof-of-concept targets, and feature breakdown. The separation into proof-of-concept targets, core features, major features, and stretch features is well organized and easy to follow. However, some proof-of-concept metrics (such as $\geq95\%$ trace mapping and $\geq80\%$ usability success rate) could briefly explain how these metrics will be measured or evaluated. This would strengthen the feasibility and evaluation clarity. Additionally, the distinction between core features and major features could be clarified slightly, as both appear required for the full system."
\end{quotation}

\begin{quotation} [20]
  "[The Architecture] section provides an excellent high-level architectural overview using the three-layer model (Ingestion, Transformation, Presentation). The layered design is clear, logical, and appropriate for the system. However, the architecture diagram should be explicitly referenced and described more clearly in the text. Explaining what each arrow and component represents would improve clarity. Additionally, some implementation details (such as exact timing offsets and layout algorithms) may be overly detailed for a proposal and could be slightly condensed."
\end{quotation}

\begin{quotation} [21]
  "[The Ingestion Layer] section clearly explains how LLVM IR is parsed and structured. The use of llvmlite and debug metadata is appropriate and well explained. However, the output format could be explained more clearly, particularly how it connects to the Transformation layer."
\end{quotation}

\begin{quotation} [22]
  "The transformation layer is well explained and provides a clear mapping between LLVM constructs and animation objects. However, the scene graph structure could benefit from a brief visual example or mock representation to clarify how data flows through the system. Additionally, the modified Sugiyama layout algorithm is mentioned but not briefly explained, which may confuse readers unfamiliar with graph layout techniques."
\end{quotation}

\begin{quotation} [23]
  "[The Presentation Layer] section clearly explains how Manim is used to render animations from the scene graph. The customization options and rendering pipeline are well described. However, the relationship between scene graph generation and Manim output could be explained more clearly in terms of data flow. Additionally, it could briefly explain expected output formats and user workflows."
\end{quotation}

\begin{quotation} [24]
  "The command-line interface examples are very helpful and clearly demonstrate usability. However, this section could briefly explain error handling and limitations more clearly, such as unsupported language features or missing debug metadata. Additionally, this section could clarify whether the CLI is intended for developers, educators, or both."
\end{quotation}

\begin{quotation} [25]
  "The projected timeline is clear and logically structured. The weekly milestones align well with the system architecture and development process. However, testing and validation phases could be more clearly defined, particularly usability testing and feature validation. Additionally, stretch feature implementation is only briefly mentioned and could be explicitly identified in the timeline."
\end{quotation}

\begin{quotation} [26]
  "The feedback section properly acknowledges prior feedback and clearly states what remains unaddressed. However, this section could be improved by explicitly explaining how each feedback item will be addressed moving forward rather than only stating that it remains unaddressed. Additionally, creating mockups or example animation output would significantly strengthen the proposal and demonstrate feasibility."
\end{quotation}

\begin{quotation} [27]
  "The references are properly formatted and relevant to the proposal. The use of academic sources and official LLVM documentation strengthens credibility. However, ensure consistent formatting across all references and verify citation formatting matches IEEE or required course format exactly."
\end{quotation}

\noindent\hrulefill

\begin{quotation} [28]
  "It's unnecessary to explain how Julia works. It distracts from the main focus of the paper."
\end{quotation}

\begin{quotation} [29]
  "You don't have to include all these accomplishments. We already get the idea that [3Blue1Brown] is popular and well respected from earlier."
\end{quotation}

\begin{quotation} [30]
  "Including pictures of Manim diagrams here would illustrate your point better."
\end{quotation}

\begin{quotation} [31]
  "Your goals seem overly ambitious for just a one semester student project."
\end{quotation}

\begin{quotation} [32]
  "It's unrealistic to expect work to be done over spring break."
\end{quotation}

\begin{quotation} [33]
  "You should talk more about the potential risks. It's not really mentioned at all."
\end{quotation}

\begin{quotation} [34]
  "The motivation section could be strengthened. While the proposal explains that existing tools are either too technical or not presentation-friendly, it does not fully explain why this problem is important at a broader level. It would help to more clearly identify who specifically benefits from this tool and why current alternatives are insufficient. The proposal currently reads as though the improvement is mostly aesthetic, so emphasizing the educational or practical impact would make the motivation more compelling."
\end{quotation}

\begin{quotation} [35]
  "The objectives are generally clear, but some could be more measurable. For example, the proposal mentions performance and trace mapping goals, but it would be helpful to specify what constitutes success more concretely. Adding clearer success criteria would make it easier to evaluate whether the project meets its goals."
\end{quotation}

\begin{quotation} [36]
  "The architectural breakdown into ingestion, transformation, and presentation layers is well structured and shows strong planning. However, the scope could be clarified further. It would be helpful to specify what subset of LLVM IR will be fully supported and how unsupported constructs will be handled. Addressing potential edge cases would make the approach feel more complete and realistic."
\end{quotation}

\begin{quotation} [37]
  "The proposal would benefit from a clearer description of final deliverables. While features are described, it is not entirely clear what concrete artifacts will be submitted at the end of the project. Explicitly listing tangible deliverables would strengthen this section."
\end{quotation}

\begin{quotation} [38]
  "The timeline is present but somewhat broad in places. Some weekly goals are described generally, which makes it difficult to determine what measurable progress will be achieved each week. Making each milestone more concrete and demonstrable would better align with the rubric’s expectations.
\end{quotation}

\begin{quotation} [39]
  "The proposal mentions evaluation goals such as trace mapping percentages and usability testing, which is good. However, the evaluation methodology could be described in more detail. For example, specifying the number of participants in the pilot study, the types of tasks they will perform, and how results will be measured would make the evaluation plan stronger and more rigorous."
\end{quotation}

\begin{quotation} [40]
  "A clearly labeled risks section would improve the proposal. While potential challenges are implied, explicitly identifying technical and timeline risks and describing mitigation strategies would demonstrate thoughtful project planning."
\end{quotation}

\noindent\hrulefill

\begin{quotation} [41]
  "I feel [the PythonTutor section] should include the supported languages of PythonTutor, since the usefulness of the tool depends heavily on what languages are supported. It is more useful if it supports languages that are commonly used for education vs if it supports languages that are more obscure."
\end{quotation}

\begin{quotation} [42]
  "What are DOT files? does this mean unix dot files, or does this stand for something"
\end{quotation}

\begin{quotation} [43]
  "What does verbosity mean in this application and how is it different from animation speed. Does the creator of the animation include an audio track, or is it the speed that words appear? If the latter, how and why are words appearing?"
\end{quotation}

\begin{quotation} [44]
  "Go more in depth about what the study entails. Will you be comparing against other similar tools, or just compared to only reading the code? Will it be between subjects or within subjects? Are you only going to collect data about accuracy or will you also collect data about subjective user preference?
  
  Also this study does not seem like it connects well to the stated goal of this project. You are discussing making a tool that will help educational scenarios, but are only evaluating it for use in debugging. I would add a component focusing on user experience generating the images, or evaluating which ones help people to better understand the purpose of code

  There are also no sample graphs or tables for the data that is to be collected"
\end{quotation}

\begin{quotation} [45]
  "Why separate core features from MVP features when core features are defined as "required to have a usable product at all", which is essentially the MVP"
\end{quotation}

\begin{quotation} [47]
  "How are animations utilized. Is it moving arrows to demonstrate code flow? Will sections of the code be moved on screen only when relevant, and removed afterward. For someone unfamiliar with Manim this is unclear"
\end{quotation}

\begin{quotation} [48]
  "While I understand that this document is written with the assumption that the reader has a computer science bachelors degree, I think it would help to use the long form of this abbreviation at least once, as was done with LLVM IR."
\end{quotation}

\begin{quotation} [49]
  "Why make your own "JSON-like" format instead of just using standard JSON. If it is for data streaming purposes, perhaps consider JSON Lines"
\end{quotation}

\begin{quotation} [50]
  "No discussion of risks or costs."
\end{quotation}

\begin{quotation} [51]
  "Timeline does not say when the experiment will be performed"
\end{quotation}

\begin{quotation} [52]
  "Missing time spent on project"
\end{quotation}

\begin{quotation} [53]
  "The lack of Manim examples or proposed examples for this project makes it difficult for someone unfamiliar with that tool to understand what this tool will produce"
\end{quotation}

% \TODO{critical} describe any feedback you did NOT take action on, with a brief justification for why (for example, if some other project related to your proposal was suggested, but it is not, this is where you would say why it was not related). Include ALL feedback you have received on the project at any point, whether from the instructor or your peers. If you have addressed every piece of feedback and made changes in response, write “We have addressed all feedback.”
